%!TEX root = ../dissertation.tex
Object-centric processes are a paradigm which represents process executions as sequences of events wherein each event may be related to different objects of multiple types, such as orders or items. Predictive process monitoring on object-centric processes seeks to provide stakeholders with actionable predictions on currently active processes by analyzing the information found in object-centric event logs. The need for more accurate predictions on a wider variety of applications has led these methodologies to adopt deep learning methods, such as neural networks, for their analysis. However, these black box methods are known for their lack of transparency, which can make their adoption in critical business areas more difficult as stakeholders are hesitant to trust models they do not properly understand.

In order to address this issue a four-component explainability framework is proposed for use with predictive process monitoring on object-centric event logs. The approach combines both factual and counterfactual explanation methods to provide stakeholders with a holistic understanding of a given model's predictions.

The framework was tested on four different regression models, trained on information from two different object-centric event logs. The result is the identification of the most impactful elements of an input on a given model's prediction and an estimation of the impact those elements have on the model's output. Additionally, through a counterfactual structure the framework is able to identify similar process executions that can help the stakeholder understand the minimum number of changes that may lead to a different prediction. These results are presented to stakeholders in the form of graphs and tables, which aim to make the information as approachable as possible to the end user.